%% parts/p05/models-and-interpretability.tex --- Phase-4 P5 depth visit
%% Source migration: v1.1 lines 9192--9229 and 9249--9260.
%% New depth: D-P6-002/003/005/006 and N-a-029/032 (partial as recorded).
%% The source exercise at 9240--9247 is retired because the handbook has no
%% exercise layer; its elementary-equivalence warning is retained below.

\chapter{Models of Arithmetic and Interpretability}
\label{sec:p05-models-and-interpretability}
\chaptermark{Models and Interpretability}

\section{Standard and nonstandard arithmetic}

\begin{hbdefinition}{def:p05-standard-nonstandard-model}{Standard and nonstandard models of arithmetic}
Let \(T\) be an \(\Lang_A\)-theory with
\(\mathcal N_{\mathrm{ar}}\entails T\).  A model
\(\mathcal M\entails T\) is \emph{standard} iff
\(\mathcal M\cong\mathcal N_{\mathrm{ar}}\) as an
\(\Lang_A\)-structure, and \emph{nonstandard} otherwise.  The elements
\(\bar n^{\mathcal M}\), for externally standard \(n\in\N\), are the
\emph{standard numerals} of \(\mathcal M\).  When
\(\mathcal M\entails\PA\), they form an order-initial copy of
\((\N,<)\), called the \emph{standard initial segment}; it is a
metalinguistically specified subset, not a predicate added to
\(\Lang_A\).  The terms standard and nonstandard classify arithmetic
structures, not sentences or proof predicates, and elementary equivalence
alone does not imply standardness.  The initial-segment terminology is used
below only for models of \(\PA\).
\end{hbdefinition}

\begin{hbtheorem}{thm:p05-compactness-nonstandard}{proved=annex:p5.models.1}{Compactness produces nonstandard models}
Assume the Axiom of Choice in the ambient set theory.  Every
\(\Lang_A\)-theory \(T\) satisfied by
\(\mathcal N_{\mathrm{ar}}\) has a countable nonstandard model.
\end{hbtheorem}

\begin{hbexample}{def:p05-nonstandard-pa}{A countable nonstandard model of PA}
Taking \(T=\PA\) in
Theorem~\ref{thm:p05-compactness-nonstandard} yields a countable model
of \(\PA\) containing an element unequal to every numeral.  The theorem
asserts existence; it neither names such an element by a standard numeral
nor makes the model isomorphic to \(\mathcal N_{\mathrm{ar}}\).
\end{hbexample}

\begin{hbdefinition}{def:p05-finite-distance-block}{Finite-distance blocks}
Let \(\mathcal M\entails\PA\) be nonstandard.  Declare two nonstandard
elements \(a,b\in\dom\mathcal M\) to lie in the same block iff, for some
externally standard \(n\in\N\),
\[
  \mathcal M\entails a+\bar n=b
  \qquad\text{or}\qquad
  \mathcal M\entails b+\bar n=a.
\]
A \emph{finite-distance block} is an equivalence class for this relation on
the nonstandard part of \(\mathcal M\).  Thus the term names an
external class of model elements, not an internally definable predicate.
It groups elements separated by finitely many externally counted successor
steps, rather than identifying elements or adjoining them to the standard
initial segment.  This distinction is needed below because within-block and
between-block order contribute different factors to the order type.
\end{hbdefinition}

\begin{hbtheorem}{thm:p05-countable-pa-order-type}{cited=Kaye1991}{Order type of a countable nonstandard model of PA}
Let \(\mathcal M\entails\PA\) be countable and nonstandard.  Under the
arithmetic order, its standard initial segment has order type \(\omega\),
each finite-distance block has order type
\(\operatorname{otp}(\mathbb Z,<)\), and the quotient order of the nonstandard
blocks is a countable dense linear order without endpoints.  Consequently
\[
  \operatorname{otp}(\dom\mathcal M,<^{\mathcal M})
  =\omega+\sum_{q\in(\mathbb Q,<)}\operatorname{otp}(\mathbb Z,<).
\]
The sum is the lexicographic sum indexed by \((\mathbb Q,<)\), with one copy of
\((\mathbb Z,<)\) at every index; it fixes the meaning of the
convention-sensitive shorthand \(\omega+\mathbb Z\cdot\mathbb Q\).  For an
uncountable nonstandard model the
block quotient remains dense without endpoints but need not have order
type \(\mathbb Q\).
\end{hbtheorem}

\begin{hbtheorem}{thm:p05-overspill}{cited=BoolosBurgessJeffrey2007}{Overspill}
Let \(\mathcal M\entails\PA\) be nonstandard, let
\(\varphi(x,\vec a)\) be an internal first-order \(\Lang_A\)-formula,
and let \(\vec a\in(\dom\mathcal M)^k\).  If
\(\mathcal M\entails\varphi(\bar n,\vec a)\) for every externally
standard \(n\in\N\), then some nonstandard \(b\in\dom\mathcal M\)
satisfies \(\mathcal M\entails\varphi(b,\vec a)\).  In the stronger
initial-segment form, some nonstandard \(c\) satisfies
\[
  \mathcal M\entails\forall x<c\,\varphi(x,\vec a).
\]
Countability is not assumed.  The conclusion concerns an internally
definable property, not an arbitrary external subset of the model
\cite{Kaye1991}.
\end{hbtheorem}

\begin{hbexample}{def:p05-overspill-below}{A nonstandard element below another}
Fix nonstandard \(c\in\dom\mathcal M\).  Every standard numeral is below
\(c\), so overspill applied to \(x<c\) gives a nonstandard \(b<c\).
Thus the standard initial segment has no immediate nonstandard successor.
Moreover, no internal formula, even with parameters from \(\mathcal M\),
defines exactly the standard initial segment: if it did, its truth at every
standard numeral would overspill to a nonstandard element.
\end{hbexample}

\begin{hbtheorem}{thm:p05-tennenbaum}{cited=Tennenbaum1959}{Tennenbaum's theorem}
No countable nonstandard \(\mathcal M\entails\PA\) has a computable
presentation.  Explicitly, there is no bijection
\(e:\N\to\dom\mathcal M\) for which both transported operations
\[
  (m,n)\longmapsto e^{-1}\!\bigl(e(m)+^{\mathcal M}e(n)\bigr),\qquad
  (m,n)\longmapsto e^{-1}\!\bigl(e(m)\times^{\mathcal M}e(n)\bigr)
\]
are computable functions on \(\N^2\).  This does not deny the existence
of countable nonstandard models, and it does not claim that their order
has no computable copy \cite{Kaye1991}.
\end{hbtheorem}

\section{Provable Sigma-one completeness}

\begin{hbtheorem}{thm:p05-sigma-one-completeness}{sketched=annex:p5.sigma.1}{External and formalized Sigma-one completeness}
Let \(\sigma\) be a closed \(\Sigma^0_1\) sentence.
\begin{enumerate}
  \item If \(\mathcal N_{\mathrm{ar}}\entails\sigma\), then
        \(\mathsf Q\derives\sigma\).  Hence every theory extending
        \(\mathsf Q\) proves every true closed \(\Sigma^0_1\) sentence.
  \item Relative to the fixed PA proof predicate of
        Definition~\ref{def:p05-proof-predicate}, the following is a
        metatheoretic schema over closed \(\Sigma^0_1\) sentences:
        \[
          \PA\derives
          \sigma\to
          \mathrm{Prov}_{\PA}(\overline{\ulcorner\sigma\urcorner}).
        \]
\end{enumerate}
The first clause is external truth-to-provability; the second internalizes
for each fixed sentence a primitive-recursive construction of a proof code
from witnesses for its bounded matrix.  Neither clause is a reflection
principle, and neither licenses
\(\mathrm{Prov}_{\PA}(\overline{\ulcorner\sigma\urcorner})\to\sigma\)
\cite{Shoenfield1967,HilbertBernays1939}.
\end{hbtheorem}

\begin{hbexample}{def:p05-sigma-one-witness}{Using a standard witness}
The true sentence \(\exists x\,(x\times x=\bar4)\) is
\(\Sigma^0_1\).  The standard witness \(2\), the numeral equations of
\(\mathsf Q\), and existential introduction give its
\(\mathsf Q\)-derivation.  The general theorem replaces this single
calculation by induction over an arbitrary bounded matrix.
\end{hbexample}

\begin{hbwarning}{def:p05-completeness-soundness}{Completeness is not soundness}
\(\Sigma^0_1\)-completeness says that true \(\Sigma^0_1\) sentences are
provable in \(\mathsf Q\).  \(\Sigma^0_1\)-soundness says that
\(\Sigma^0_1\) sentences proved by a theory are true in
\(\mathcal N_{\mathrm{ar}}\).  The implications point in opposite
directions and are used at different steps of the incompleteness arguments.
\end{hbwarning}

\section{Relative interpretation}

\begin{hbdefinition}{def:p05-relative-interpretation}{One-dimensional parameter-free relative interpretation}
Let \(U\) and \(V\) be theories in signatures \(\Lang_U\) and
\(\Lang_V\).  A \emph{one-dimensional parameter-free interpretation}
\(\mathcal I\) of \(U\) in \(V\) consists of a \(\Lang_V\)-formula
\(\varphi_{\mathrm{dom}}(x)\), with
\(V\derives\exists x\,\varphi_{\mathrm{dom}}(x)\), an interpreted-equality
formula \(E(x,y)\), and the following data.  One may take \(E(x,y)\) to be
target equality; in the quotient case, \(V\) instead proves that \(E\) is
an equivalence relation on the interpreted domain.
\begin{enumerate}
  \item a \(\Lang_V\)-formula for every \(\Lang_U\)-relation, invariant
        under \(E\) in each argument;
  \item a \(\Lang_V\)-definable graph \(F_f(\vec x,y)\) for every
        \(\Lang_U\)-function \(f\), treating constants as nullary
        functions.  The theory \(V\) proves domain closure and totality,
        uniqueness of outputs modulo \(E\), and invariance under \(E\) in
        every input and the output;
  \item a recursive translation \(\psi\mapsto\psi^{\mathcal I}\) that
        translates source equality as \(E\), is induced on atomic formulas
        by the relation and graph formulas, commutes with Boolean
        connectives, and replaces quantifiers by
        \[
          (\exists x\,\psi)^{\mathcal I}\defeq
          \exists x\,(\varphi_{\mathrm{dom}}(x)\land\psi^{\mathcal I}),
          \qquad
          (\forall x\,\psi)^{\mathcal I}\defeq
          \forall x\,(\varphi_{\mathrm{dom}}(x)\to\psi^{\mathcal I});
        \]
  \item \(V\derives\alpha^{\mathcal I}\) for every axiom
        \(\alpha\) of \(U\).
\end{enumerate}
Here uniqueness modulo \(E\) means that two outputs of the same graph on
\(E\)-equivalent inputs are \(E\)-equivalent; this avoids silently assuming
a definable choice of quotient representatives.  The term names
metalinguistically specified translation data between theories, not an
arithmetic object or a model.  It is needed because literal inclusion and
definitional extension alone cannot express a definable subdomain or
quotient inside models of \(V\).  It is stronger data than a mere
consistency comparison and weaker than definitional equivalence or
bi-interpretability, and it is used in the transfer theorem below
\cite{TMR1953,Lindstrom2003}.
\end{hbdefinition}

\begin{hbdefinition}{def:p05-essential-undecidability}{Essential undecidability}
A consistent theory \(U\) is \emph{essentially undecidable} iff every
consistent computably enumerable theory in the same signature that
extends \(U\) is undecidable.  This is a metatheoretic property of a
theory.  It is stronger than undecidability of \(U\) itself because it
also constrains all consistent effective extensions, and that stronger
distinction is required in the transfer theorem.
\end{hbdefinition}

\begin{hbtheorem}{thm:p05-interpretation-transfer}{cited=TMR1953}{Interpretation transfer}
If \(\mathcal I\) interprets \(U\) in \(V\), then
\[
  U\derives\psi\quad\Longrightarrow\quad
  V\derives\psi^{\mathcal I}.
\]
Consequently, consistency of \(V\) implies consistency of \(U\).  If,
in addition, \(U\) is essentially undecidable, \(V\) is consistent and
computably enumerable, and the interpretation is effective, then \(V\)
is essentially undecidable and hence undecidable.  Plain undecidability
alone does not generally transfer under an arbitrary interpretation.
\end{hbtheorem}

\begin{hbexample}{def:p05-interpretation-examples}{Two elementary interpretations}
Because \(\PA\) extends \(\mathsf Q\), the identity translation interprets
\(\mathsf Q\) in \(\PA\).  A definitional extension interprets back into
its base theory by replacing each new symbol by its defining formula;
unlike the identity case, this second translation removes vocabulary.
These examples are interpretations, but the definition does not make
arbitrary interpretations conservative or invertible.
\end{hbexample}

Interpretability logic studies modal principles valid under arithmetical
interpretability.  Its calculi and completeness theory are outside this
visit; see \cite{Visser1997} and the Phase-5 independence chapter.
